\chapter{Metodo sperimentale}
\label{cap:22-metodo-sperimentale}

I capitoli seguenti riportano misure. Questo descrive come sono state ottenute,
e --- pi\`u importante --- quali propriet\`a dell'impianto rendono confrontabili
due numeri raccolti in momenti diversi.

\section{Il principio: comparabilit\`a nel tempo, non fra piattaforme}

Il \cref{cap:02-stato-dell-arte} ha osservato che confrontare piattaforme FaaS
misurate da persone diverse su hardware diverso \`e notoriamente inaffidabile.
Il metodo di \nf{} non tenta quel confronto. Persegue invece la comparabilit\`a
\emph{intra}-piattaforma: fissati un profilo hardware e un carico, il risultato
di una release \`e confrontato con quello della release precedente sullo stesso
profilo, e la deviazione \`e trattata come segnale di regressione.

\`E un obiettivo pi\`u ristretto e pi\`u difendibile. Non consente di affermare
che \nf{} sia pi\`u veloce di Knative; consente di affermare che una certa
modifica ha cambiato il throughput del $x\%$ su un profilo dichiarato, che \`e
l'affermazione di cui il progetto ha effettivamente bisogno.

\section{L'impianto: uno strumento separato}

L'orchestrazione degli esperimenti vive in un repository separato. La
separazione \`e deliberata e ha una giustificazione di metodo: uno strumento che
provvisiona macchine, installa Kubernetes, distribuisce immagini ed esegue
carichi \`e un sistema a s\'e, con dipendenze proprie, e includerlo nel
repository della piattaforma renderebbe la piattaforma pi\`u difficile da
costruire per chi vuole solo usarla.

Il confine \`e quello enunciato nel \cref{cap:19-cli-scaffolding}: la CLI \`e un
client HTTP e non provvisiona nulla; lo strumento di provisioning non conosce il
protocollo delle funzioni.

\subsection{Scenari dichiarativi}

Un esperimento \`e descritto da uno scenario YAML e da un file di ambiente. Lo
scenario dichiara \emph{che cosa} accade --- provvisiona, installa, registra,
carica, verifica --- e l'ambiente dichiara \emph{dove}.

\begin{lstlisting}[language=bash]
# esecuzione
./nanolab.sh run packages/nanolab/scenarios-v2/deployment-lifecycle-k8s.yaml \
    --environment packages/nanolab/environments/multipass.yaml
# anteprima senza eseguire
./nanolab.sh plan packages/nanolab/scenarios-v2/deployment-lifecycle-k8s.yaml \
    --environment packages/nanolab/environments/local.yaml
\end{lstlisting}

La validazione end-to-end della piattaforma appartiene a questo strumento e non
al repository di \nf{}: la suite di quest'ultimo si ferma ai test unitari e di
integrazione, e il ciclo di vita completo --- registrazione, invocazione,
verifiche HTTP e Kubernetes --- \`e uno scenario
(\cref{cap:26-qualita-verifica}). Le famiglie di scenari sono declinate per
backend di esecuzione: lo stesso ciclo di vita esiste per il backend container,
per containerd e per Kubernetes (\code{deployment-lifecycle-container},
\code{-containerd}, \code{-k8s}), cos\`i che un cambiamento di backend non
richieda di riscrivere l'esperimento; al 2026-09-28 la cartella
\code{scenarios-v2} contiene 70 file.

L'esistenza di una modalit\`a \code{plan} \`e significativa: un esperimento che
provvisiona infrastruttura a pagamento non dovrebbe scoprire un errore di
configurazione dopo dieci minuti di provisioning.

\subsection{Ambienti disponibili}

\begin{center}
\small
\begin{adjustbox}{max width=\textwidth}
\begin{tabular}{@{}lp{7.6cm}@{}}
\toprule
\textbf{Ambiente} & \textbf{Uso} \\
\midrule
locale (Docker) & Validazione del container, cicli rapidi. \\
Multipass & VM locale con k3s: percorso Kubernetes completo senza costo. \\
Azure & Profilo fissato per le misure pubblicate; VM dedicate per stack e
generatore di carico. \\
Proxmox & Ambiente su hardware proprio. \\
\bottomrule
\end{tabular}
\end{adjustbox}
\end{center}

\section{Il generatore di carico}

Il carico \`e generato con k6~\cite{k6}. Due propriet\`a dell'impianto meritano
attenzione perch\'e hanno influenzato i risultati.

\subsection{Corpora di payload versionati}

L'input dei benchmark non \`e generato casualmente ma letto da corpora
versionati (\code{functions/test-data/<famiglia>/performance-<profilo>.json}),
con tre profili di dimensione. Due esecuzioni della stessa configurazione
ricevono quindi esattamente lo stesso input.

\`E la stessa distinzione discussa nel \cref{cap:18-sdk}: i corpora di
prestazione omettono deliberatamente gli output attesi, perch\'e verificarli
durante un benchmark aggiungerebbe alla misura il costo del confronto.

\subsection{Ciclo chiuso e ciclo aperto}

Questa \`e la distinzione metodologicamente pi\`u importante del capitolo, e la
sua rilevanza \`e stata scoperta a caro prezzo.

In un generatore a \textbf{ciclo chiuso}, ogni utente virtuale mantiene una sola
richiesta in volo e ne invia un'altra solo dopo aver ricevuto la risposta. Il
numero di richieste nel sistema \`e quindi fissato dal generatore. In un
generatore a \textbf{ciclo aperto}, gli arrivi sono programmati dall'orologio
indipendentemente da quanto il sistema stia impiegando a rispondere.

La differenza \`e cruciale per valutare un controllore di concorrenza:

\begin{quote}\itshape
Sotto il profilo a raffica ogni utente virtuale tiene una richiesta, quindi il
numero nel sistema \`e fissato dal generatore e la profondit\`a di coda \`e
\code{VU} $-$ \code{limite}. Due conseguenze, entrambe misurate. Portare il
limite da 4 a 8 ha accorciato di quattro una coda profonda 99 --- nessun
controllore pu\`o muovere l'attesa di molto. E la coda si assesta a 95 su 100 per
costruzione, permanentemente sopra qualunque soglia di guardia, quindi la
guardia di ammissione scatta a ogni tick e il modello non parla mai.
\end{quote}

La conseguenza \`e enunciata senza mezzi termini: \emph{``\`E per questo che due
controllori che leggevano segnali completamente diversi sono finiti entro il 2\%
l'uno dall'altro su questo profilo. La somiglianza era una propriet\`a
dell'impianto.''}

\`E un risultato metodologico che va oltre questo progetto. Un carico a ciclo
chiuso non pu\`o porre al sistema la domanda a cui un controllore di ammissione
dovrebbe rispondere, perch\'e retroaziona sulla stessa grandezza che il
controllore manipola. Il \cref{cap:25-valutazione-concorrenza} riporta i due
insiemi di risultati, con e senza questa correzione.

\section{Profili hardware fissati}

Ogni record di prestazione porta l'identificatore del profilo su cui \`e stato
ottenuto:

\begin{lstlisting}[language=bash]
azure-d8s-v5+d2s-v5-amd64-native-loadtest-v1
\end{lstlisting}

L'identificatore codifica: provider, tipo di VM dello stack, tipo di VM del
generatore, architettura, tipo di build del control plane, e una versione del
profilo stesso. Un confronto fra record di profili diversi \`e vietato per
costruzione.

Il generatore di carico gira su una VM \emph{separata} da quella dello stack. \`E
una precauzione elementare e spesso omessa: un generatore co-residente compete
per CPU con il sistema misurato, e il suo contributo alla latenza osservata
cresce proprio quando il sistema \`e sotto pressione.

\section{Il record di release}

Ogni release produce un record JSON versionato in
\code{docs/performance/releases/}. Contiene:

\begin{center}
\small
\begin{adjustbox}{max width=\textwidth}
\begin{tabular}{@{}lp{7.6cm}@{}}
\toprule
\textbf{Sezione} & \textbf{Contenuto} \\
\midrule
\code{aggregates} & Throughput, tasso di errore, p50/p95/p99, cold start,
repliche di picco, CPU e heap di picco del control plane, attesa in coda. \\
\code{profile} & L'identificatore del profilo e le sue componenti. \\
\code{runCount} & Numero di esecuzioni la cui mediana \`e riportata (3). \\
\code{sourceCommit} & Il commit esatto da cui l'artefatto \`e stato costruito. \\
\code{imageDigests} & Digest SHA-256 di \emph{ogni} immagine coinvolta. \\
\code{thresholds} & Le soglie di regressione applicate. \\
\bottomrule
\end{tabular}
\end{adjustbox}
\end{center}

I digest sono la parte che rende il record realmente riproducibile: non
identificano un'etichetta mobile ma il contenuto esatto delle immagini. Un
record di \nf{} v0.19.0 elenca oltre centoventi digest --- il control plane in
quattro varianti, il watchdog, e le funzioni di riferimento in cinque linguaggi
per due architetture.

\section{Il gate di regressione}

Tre esecuzioni per release, mediane confrontate con il record pi\`u recente dello
stesso profilo. Le soglie:

\begin{center}
\small
\begin{adjustbox}{max width=\textwidth}
\begin{tabular}{@{}lrl@{}}
\toprule
\textbf{Soglia} & \textbf{Valore} & \textbf{Significato} \\
\midrule
\code{errorRateMax} & 0{,}3\,\% & Tasso di errore massimo assoluto. \\
\code{p95MaxIncreasePercent} & 15\,\% & Peggioramento massimo del p95 rispetto
al record precedente. \\
\code{throughputMaxLossPercent} & 10\,\% & Perdita massima di throughput. \\
\bottomrule
\end{tabular}
\end{adjustbox}
\end{center}

Le due soglie relative sono generose --- 15\% e 10\% --- e la ragione \`e onesta:
con tre esecuzioni la varianza campionaria non consente di distinguere
differenze pi\`u piccole. Una soglia stretta produrrebbe falsi allarmi a ogni
release, e il gate verrebbe disattivato. Il \cref{cap:27-discussione} riprende
il punto come limite alla sensibilit\`a del metodo.

Accanto alle soglie numeriche esiste un gate strutturale, gi\`a citato nel
\cref{cap:21-osservabilita}: una modifica che regredisca le garanzie
strutturali del percorso caldo --- riuso del replay, forward progress della coda
sincrona, equit\`a asincrona --- viene respinta anche se i numeri assoluti
sembrano accettabili.

\section{Un errore di calibrazione, riportato}

La documentazione del progetto conserva un errore che vale la pena riportare
perch\'e \`e generale:

\begin{quote}\itshape
La prima esecuzione di questo confronto ha fissato il budget a 8 mentre il carico
raggiungeva un picco di 105 contro una coda di 100. Sotto contesa ciascuna
funzione teneva circa 4 posti di 104 contro 105 arrivi, quindi la coda traboccava
per aritmetica in un modo e non nell'altro, producendo 86\,302 rifiuti che non
dicevano nulla sul controllore.
\end{quote}

E la regola che ne \`e stata derivata:

\begin{quote}\itshape
Il budget deve stare dentro una finestra: a o sopra \emph{funzioni $\times$
tetto} nulla \`e mai scarso e il modo degenera in controllo per funzione; sotto
\emph{funzioni $\times$ (picco $-$ coda)} i rifiuti sono fabbricati. Il picco
offerto deve essere derivato dal limite pi\`u piccolo che qualunque modo sotto
confronto conceder\`a, o le due esecuzioni non stanno misurando la stessa cosa.
\end{quote}

\`E un vincolo che non compare in nessuna documentazione di benchmarking
serverless, e che discende dalla struttura del sistema misurato: quando la
grandezza controllata determina anche la finestra di ammissione, il carico
offerto non pu\`o essere scelto indipendentemente dal controllore.

\section{Riproducibilit\`a: che cosa \`e garantito e che cosa no}

\begin{center}
\small
\begin{adjustbox}{max width=\textwidth}
\begin{tabular}{@{}lp{7.8cm}@{}}
\toprule
\textbf{Elemento} & \textbf{Stato} \\
\midrule
Codice sorgente & Fissato dal commit nel record. \\
Immagini & Fissate dai digest nel record. \\
Input del carico & Fissato dai corpora versionati. \\
Profilo hardware & Fissato dall'identificatore; la VM \`e ricreata a ogni
esecuzione. \\
Prestazioni della VM & \textbf{Non garantite}: due VM Azure dello stesso tipo
possono avere vicini diversi. \\
Numero di ripetizioni & 3, che \`e poco. \\
\bottomrule
\end{tabular}
\end{adjustbox}
\end{center}

Le ultime due righe sono i limiti reali. Il \cref{cap:27-discussione} li tratta
per esteso, ma \`e utile enunciarli qui, accanto al metodo che pretendono di
qualificare: le misure riportate nei prossimi tre capitoli sono \emph{mediane di
tre esecuzioni su infrastruttura condivisa}, non stime con intervallo di
confidenza.
